
Large language models exhibit hallucinations on factual questions, necessitating uncertainty quantification for deployment in domains where incorrect outputs carry consequences. Semantic entropy (Farquhar et al., 2024) computes uncertainty at the meaning level by generating multiple samples, clustering semantically equivalent responses, and calculating discrete entropy over clusters. This method achieves AUROC approximately 0.80 for hallucination detection on benchmarks but requires K=10 samples followed by natural language inference operations, resulting in latency of approximately 1500ms per query.

Production deployment of language models in latency-sensitive applications requires response times below 100ms. The 15-fold overhead imposed by semantic entropy computation creates a barrier to deploying uncertainty-aware systems in contexts where users expect sub-second interactions. Existing fast alternatives operate at the token level (perplexity, token probabilities) and correlate weakly with semantic-level uncertainty. Supervised probes (Kossen et al., 2024) achieve accuracy comparable to semantic entropy but require training on uncertainty labels and produce non-interpretable predictions.

We hypothesized that intrinsic geometric properties of hidden state manifolds could provide an alternative approach. If epistemic uncertainty manifests as lower-dimensional structure in late-layer representations, then spectral features computable from covariance matrices might correlate with semantic entropy without requiring multi-sample generation. Specifically, we proposed that participation ratio (effective dimensionality), eigenvalue decay rate, and condition number extracted from layers 24-31 at the final token position would achieve Spearman correlation magnitude exceeding 0.4 with semantic entropy on TruthfulQA factual questions.

We designed and implemented a complete experimental pipeline to test this hypothesis using 7B parameter models on 817 TruthfulQA questions (246 test examples). However, hidden state extraction proved computationally infeasible, consuming more than 10 hours without producing output for the test set. This computational failure prevented all empirical validation. Zero correlation measurements were obtained. The research question remains unanswered.

This negative result documents a computational constraint others pursuing geometric approaches to uncertainty quantification will encounter. We quantify resource requirements through systematic analysis of the failure point and specify conditions under which validation becomes tractable.
